\documentclass{article}
\usepackage{fleqn}
\usepackage{epsf}
\usepackage{aima2e-slides}

\begin{document}

\begin{huge}
\titleslide{Thị giác máy tính (Vision)}{Chương 24}

\sf

%%%%%%%%%%%% Slide %%%%%%%%%%%%%%%%%%%%%%%%%%%%%%%%%%%%%%%%%%%%%%%%%%%
\heading{Phác thảo}

\blob Nhận thức chung

\blob Sự hình thành hình ảnh

\blob Tầm nhìn sớm

\blob 2D \mat{$\rightarrow$} 3D

\blob Nhận dạng đối tượng




%%%%%%%%%%%% Slide %%%%%%%%%%%%%%%%%%%%%%%%%%%%%%%%%%%%%%%%%%%%%%%%%%%
\heading{Nhận thức chung}

\defn{Kích thích} (nhận thức) \mat{$S$}, Thế giới \mat{$W$}
\mat{\[
  S = g(W)
\]}%
Ví dụ: \mat{$g$} = ``graphics.'' Chúng ta có thể thực hiện tầm nhìn như đồ họa nghịch đảo không?
\mat{\[
  W = g^{-1}(S)
\]}%

%%%%%%%%%%%% Slide %%%%%%%%%%%%%%%%%%%%%%%%%%%%%%%%%%%%%%%%%%%%%%%%%%%
\heading{Nhận thức chung}

\defn{Kích thích} (nhận thức) \mat{$S$}, Thế giới \mat{$W$}
\mat{\[
  S = g(W)
\]}%
Ví dụ: \mat{$g$} = ``graphics.'' Chúng ta có thể thực hiện tầm nhìn như đồ họa nghịch đảo không?
\mat{\[
  W = g^{-1}(S)
\]}%
Vấn đề: sự mơ hồ lớn!

\vspace*{0.2in}

\epsfxsize=1.0\textwidth
\fig{\sfile{figures}{lecture-scene1.ps}}

%%%%%%%%%%%% Slide %%%%%%%%%%%%%%%%%%%%%%%%%%%%%%%%%%%%%%%%%%%%%%%%%%%
\heading{Nhận thức chung}

\defn{Kích thích} (nhận thức) \mat{$S$}, Thế giới \mat{$W$}
\mat{\[
  S = g(W)
\]}%
Ví dụ: \mat{$g$} = ``graphics.'' Chúng ta có thể thực hiện tầm nhìn như đồ họa nghịch đảo không?
\mat{\[
  W = g^{-1}(S)
\]}%
Vấn đề: sự mơ hồ lớn!

\vspace*{0.2in}

\epsfxsize=1.0\textwidth
\fig{\sfile{figures}{lecture-scene2.ps}}

%%%%%%%%%%%% Slide %%%%%%%%%%%%%%%%%%%%%%%%%%%%%%%%%%%%%%%%%%%%%%%%%%%
\heading{Nhận thức chung}

\defn{Kích thích} (nhận thức) \mat{$S$}, Thế giới \mat{$W$}
\mat{\[
  S = g(W)
\]}%
Ví dụ: \mat{$g$} = ``graphics.'' Chúng ta có thể thực hiện tầm nhìn như đồ họa nghịch đảo không?
\mat{\[
  W = g^{-1}(S)
\]}%
Vấn đề: sự mơ hồ lớn!

\vspace*{0.2in}

\epsfxsize=1.0\textwidth
\fig{\sfile{figures}{lecture-scene3.ps}}


%%%%%%%%%%%% Slide %%%%%%%%%%%%%%%%%%%%%%%%%%%%%%%%%%%%%%%%%%%%%%%%%%%
\heading{Phương pháp tiếp cận tốt hơn}

Suy luận Bayes về cấu hình thế giới:
\mat{\[
  P(W|S) = \alpha \underbrace{P(S|W)}_{\mbox{``graphics''}}\quad \underbrace{P(W)}_{\mbox{``prior knowledge''}}
\]}%
Vẫn còn tốt hơn: không cần phải khôi phục cảnh chính xác!\\
Chỉ cần trích xuất thông tin cần thiết cho \al
 -- \defn{điều hướng}\al
 -- \defn{thao tác}\al
 -- \defn{nhận dạng/nhận dạng}


%%%%%%%%%%%% Slide %%%%%%%%%%%%%%%%%%%%%%%%%%%%%%%%%%%%%%%%%%%%%%%%%%%
\heading{Vision ``hệ thống con''}

\vspace*{0.2in}

\epsfxsize=1.0\textwidth
\fig{\file{figures}{vision-subsystems.ps}}

Tầm nhìn đòi hỏi phải kết hợp nhiều tín hiệu

%%%%%%%%%%%% Slide %%%%%%%%%%%%%%%%%%%%%%%%%%%%%%%%%%%%%%%%%%%%%%%%%%%
\heading{Hình thành hình ảnh}

\vspace*{0.2in}

\epsfxsize=0,75\textwidth
\fig{\file{figures}{pinhole.ps}}

\mat{$P$} là một điểm trong cảnh, có tọa độ \mat{$ (X,Y,Z)$}\\
\mat{$P'$} là ảnh của nó trên mặt phẳng ảnh, có tọa độ \mat{$ (x,y,z)$}
\mat{\[x=\frac{-fX}{Z},\ y=\frac{-fY}{Z}  \]}%
bằng các tam giác đồng dạng. Tỷ lệ/khoảng cách là không xác định!
  

%%%%%%%%%%%% Slide %%%%%%%%%%%%%%%%%%%%%%%%%%%%%%%%%%%%%%%%%%%%%%%%%%%
\heading{Hình ảnh}

\vspace*{0.2in}

\epsfxsize=0,85\textwidth
\fig{\file{figures}{stapler1+square.ps}}



%%%%%%%%%%%% Slide %%%%%%%%%%%%%%%%%%%%%%%%%%%%%%%%%%%%%%%%%%%%%%%%%%%
\heading{Hình ảnh tiếp theo.}

\vspace*{0.2in}

\twofig{\file{figures}{pixels-12x12.ps}}{\file{figures}{pixel-values.ps}}

\mat{$I(x,y,t)$} là cường độ tại \mat{$(x,y)$} tại thời điểm \mat{$t$}

Máy ảnh CCD \mat{$\approx$} 1.000.000 pixel; mắt người \mat{$\approx$} 240.000.000 pixel\\
tức là 0,25 terabit/giây

%%%%%%%%%%%% Slide %%%%%%%%%%%%%%%%%%%%%%%%%%%%%%%%%%%%%%%%%%%%%%%%%%%
\heading{Tầm nhìn màu sắc}

Cường độ thay đổi theo tần số \mat{$\rightarrow$} tín hiệu vô chiều

\vspace*{0.2in}

\epsfxsize=0,65\textwidth
\fig{\file{figures}{color-vision.ps}}

Mắt người có ba loại tế bào nhạy cảm với màu sắc;\\
mỗi tín hiệu tích hợp cường độ vectơ 3 phần tử \mat{$\implies$}



%%%%%%%%%%%% Slide %%%%%%%%%%%%%%%%%%%%%%%%%%%%%%%%%%%%%%%%%%%%%%%%%%%
\heading{Phát hiện cạnh}

\vspace*{0.2in}

\epsfxsize=0,7\textwidth
\fig{\file{figures}{stapler1.edge-test.ps}}

\vspace*{-1in}

Các cạnh trong ảnh \mat{$\Leftarrow$} sự gián đoạn trong cảnh:\al
1) độ sâu\al
2) hướng bề mặt\al
3) độ phản xạ (dấu hiệu bề mặt)\al
4) chiếu sáng (bóng tối, v.v.)


%%%%%%%%%%%% Slide %%%%%%%%%%%%%%%%%%%%%%%%%%%%%%%%%%%%%%%%%%%%%%%%%%%
\heading{Tiếp tục phát hiện cạnh}

1) Kết hợp hình ảnh với các bộ lọc định hướng không gian (có thể đa tỷ lệ)
\mat{\[
  E_{\theta}(x,y) = \int_{-\infty}^{\infty} \int_{-\infty}^{\infty} f_{\theta}(u,v) I(x+u,y+v)\,du\,dv
\]}%

\vspace*{-0.3in}

\epsfxsize=0,75\textwidth
\fig{\file{figures}{spatially-oriented-filters.ps}}

\vspace*{-0.3in}

2) Gắn nhãn các pixel trên ngưỡng với hướng cạnh

3) Suy ra các đoạn đường ``sạch'' bằng cách kết hợp các pixel cạnh có cùng hướng

\vspace*{0.2in}

\epsfxsize=0,5\textwidth
\fig{\file{figures}{edgels.ps}}

  


%%%%%%%%%%%% Slide %%%%%%%%%%%%%%%%%%%%%%%%%%%%%%%%%%%%%%%%%%%%%%%%%%%
\heading{Tín hiệu từ kiến thức trước đây}

\begin{tabular}{l|l}
\tabbot {\bf Hình dạng từ\mat{$\ldots$} & {\bf Giả sử} \\
\hline
\chuyển động đỉnh và vật rắn, chuyển động liên tục \\
âm thanh nổi & khối liền kề, không lặp lại \\
kết cấu & kết cấu đồng đều \\
bóng và phản xạ đồng đều \\
đường viền & độ cong tối thiểu
\end{tabular}


%%%%%%%%%%%% Slide %%%%%%%%%%%%%%%%%%%%%%%%%%%%%%%%%%%%%%%%%%%%%%%%%%%
\heading{Chuyển động}

\noindent

\epsffile{\file{figures}{Rubic.00.ps}}

\centerline{%
\epsfxsize=0,3\textwidth
\noindent \epsffile{\file{figures}{Rubic.00.ps}} %
\vspace*{0.1in}\fig{\file{figures}{Rubic.flow.ps}}
\epsfxsize=0,3\textwidth
\epsffile{\file{figures}{Rubic.19.ps}}}

\vspace*{0.1in}

\epsfxsize=0,4\textwidth
\fig{\file{figures}{Rubic.flow.ps}}


%%%%%%%%%%%% Slide %%%%%%%%%%%%%%%%%%%%%%%%%%%%%%%%%%%%%%%%%%%%%%%%%%%
\heading{Âm thanh nổi}

\vspace*{0.2in}

\epsfxsize=0,85\textwidth
\fig{\file{figures}{stereo-pyramid.ps}}

%%%%%%%%%%%% Slide %%%%%%%%%%%%%%%%%%%%%%%%%%%%%%%%%%%%%%%%%%%%%%%%%%%
\heading{Độ phân giải độ sâu âm thanh nổi }

\vspace*{0.2in}

\epsfxsize=0.9\textwidth
\fig{\file{figures}{stereopsis.ps}}

Hình học đơn giản: \mat{$\delta Z = Z^2 \delta \theta/(-b)$}\\
Sinh lý học: \mat{$\delta\theta \geq 2.42\times 10^{-5}$} radian, \mat{$b\eq 6$}cm

\mat{$Z\eq 30$}cm \mat{$\implies \delta Z \approx 0.04$}mm\\
\mat{$Z\eq 30$}m \mat{$\implies \delta Z \approx 40$}cm

Đường cơ sở lớn \mat{$\implies$} có độ phân giải tốt hơn!



%%%%%%%%%%%% Slide %%%%%%%%%%%%%%%%%%%%%%%%%%%%%%%%%%%%%%%%%%%%%%%%%%%
\heading{Hoạ tiết}

\vspace*{0.2in}

\epsfxsize=0,6\textwidth
\fig{\file{figures}{chem-test.ps}}

Ý tưởng: giả sử kết cấu thực tế là đồng nhất, tính toán hình dạng bề mặt sẽ tạo ra sự biến dạng này

Ý tưởng tương tự cũng áp dụng cho việc tạo bóng---giả sử hệ số phản xạ đồng đều, v.v.---\emph{nhưng}\\
sự giao thoa đưa ra tính toán phi tiêu điểm về cường độ cảm nhận \\
\mat{$\implies$} những chỗ trũng có vẻ nông hơn thực tế



%%%%%%%%%%%% Slide %%%%%%%%%%%%%%%%%%%%%%%%%%%%%%%%%%%%%%%%%%%%%%%%%%%
\heading{Các loại cạnh và đỉnh}

\vspace*{0.2in}

\epsfxsize=0,6\textwidth
\fig{\file{figures}{diff-labels.ps}}

\epsfxsize=0,6\textwidth
\fig{\file{figures}{trihedral.ps}}

Giả sử thế giới của các vật thể đa diện rắn có đỉnh tam diện


%%%%%%%%%%%% Slide %%%%%%%%%%%%%%%%%%%%%%%%%%%%%%%%%%%%%%%%%%%%%%%%%%%
\heading{Nhãn đỉnh/cạnh}

\vspace*{0.2in}

\epsfxsize=0,6\textwidth
\fig{\file{figures}{vertex-a.ps}}

\vspace*{0.2in}

\epsfxsize=0,5\textwidth
\fig{\file{figures}{huffman-clowes.ps}}


%% %%%%%%%%%%%% Slide %%%%%%%%%%%%%%%%%%%%%%%%%%%%%%%%%%%%%%%%%%%%%%%%%%%
%% \heading{Vertex/edge labelling example}
%% 
%% \vspace*{0.2in}
%% 
%% \twofig{\sfile{figures}{edge-labelling01.ps}}{\file{figures}{huffman-clowes.ps}}
%% 
%% %%%%%%%%%%%% Slide %%%%%%%%%%%%%%%%%%%%%%%%%%%%%%%%%%%%%%%%%%%%%%%%%%%
%% \heading{Vertex/edge labelling example}
%% 
%% \vspace*{0.2in}
%% 
%% \twofig{\sfile{figures}{edge-labelling02.ps}}{\file{figures}{huffman-clowes.ps}}
%% 
%% %%%%%%%%%%%% Slide %%%%%%%%%%%%%%%%%%%%%%%%%%%%%%%%%%%%%%%%%%%%%%%%%%%
%% \heading{Vertex/edge labelling example}
%% 
%% \vspace*{0.2in}
%% 
%% \twofig{\sfile{figures}{edge-labelling03.ps}}{\file{figures}{huffman-clowes.ps}}
%% 
%% %%%%%%%%%%%% Slide %%%%%%%%%%%%%%%%%%%%%%%%%%%%%%%%%%%%%%%%%%%%%%%%%%%
%% \heading{Vertex/edge labelling example}
%% 
%% \vspace*{0.2in}
%% 
%% \twofig{\sfile{figures}{edge-labelling04.ps}}{\file{figures}{huffman-clowes.ps}}
%% 
%% %%%%%%%%%%%% Slide %%%%%%%%%%%%%%%%%%%%%%%%%%%%%%%%%%%%%%%%%%%%%%%%%%%
%% \heading{Vertex/edge labelling example}
%% 
%% \vspace*{0.2in}
%% 
%% \twofig{\sfile{figures}{edge-labelling05.ps}}{\file{figures}{huffman-clowes.ps}}
%% 
%% %%%%%%%%%%%% Slide %%%%%%%%%%%%%%%%%%%%%%%%%%%%%%%%%%%%%%%%%%%%%%%%%%%
%% \heading{Vertex/edge labelling example}
%% 
%% \vspace*{0.2in}
%% 
%% \twofig{\sfile{figures}{edge-labelling06.ps}}{\file{figures}{huffman-clowes.ps}}
%% 
%% %%%%%%%%%%%% Slide %%%%%%%%%%%%%%%%%%%%%%%%%%%%%%%%%%%%%%%%%%%%%%%%%%%
%% \heading{Vertex/edge labelling example}
%% 
%% \vspace*{0.2in}
%% 
%% \twofig{\sfile{figures}{edge-labelling07.ps}}{\file{figures}{huffman-clowes.ps}}
%% 
%% %%%%%%%%%%%% Slide %%%%%%%%%%%%%%%%%%%%%%%%%%%%%%%%%%%%%%%%%%%%%%%%%%%
%% \heading{Vertex/edge labelling example}
%% 
%% \vspace*{0.2in}
%% 
%% \twofig{\sfile{figures}{edge-labelling08.ps}}{\file{figures}{huffman-clowes.ps}}
%% 
%% %%%%%%%%%%%% Slide %%%%%%%%%%%%%%%%%%%%%%%%%%%%%%%%%%%%%%%%%%%%%%%%%%%
%% \heading{Vertex/edge labelling example}
%% 
%% \vspace*{0.2in}
%% 
%% \twofig{\sfile{figures}{edge-labelling09.ps}}{\file{figures}{huffman-clowes.ps}}
%% 
%% %%%%%%%%%%%% Slide %%%%%%%%%%%%%%%%%%%%%%%%%%%%%%%%%%%%%%%%%%%%%%%%%%%
%% \heading{Vertex/edge labelling example}
%% 
%% \vspace*{0.2in}
%% 
%% \twofig{\sfile{figures}{edge-labelling06.ps}}{\file{figures}{huffman-clowes.ps}}
%% 
%% %%%%%%%%%%%% Slide %%%%%%%%%%%%%%%%%%%%%%%%%%%%%%%%%%%%%%%%%%%%%%%%%%%
%% \heading{Vertex/edge labelling example}
%% 
%% \vspace*{0.2in}
%% 
%% \twofig{\sfile{figures}{edge-labelling10.ps}}{\file{figures}{huffman-clowes.ps}}
%% 
%% %%%%%%%%%%%% Slide %%%%%%%%%%%%%%%%%%%%%%%%%%%%%%%%%%%%%%%%%%%%%%%%%%%
%% \heading{Vertex/edge labelling example}
%% 
%% \vspace*{0.2in}
%% 
%% \twofig{\sfile{figures}{edge-labelling11.ps}}{\file{figures}{huffman-clowes.ps}}
%% 
%% %%%%%%%%%%%% Slide %%%%%%%%%%%%%%%%%%%%%%%%%%%%%%%%%%%%%%%%%%%%%%%%%%%
%% \heading{Vertex/edge labelling example}
%% 
%% \vspace*{0.2in}
%% 
%% \twofig{\sfile{figures}{edge-labelling12.ps}}{\file{figures}{huffman-clowes.ps}}
%% 
%% %%%%%%%%%%%% Slide %%%%%%%%%%%%%%%%%%%%%%%%%%%%%%%%%%%%%%%%%%%%%%%%%%%
%% \heading{Vertex/edge labelling example}
%% 
%% \vspace*{0.2in}
%% 
%% \twofig{\sfile{figures}{edge-labelling10.ps}}{\file{figures}{huffman-clowes.ps}}
%% 
%% %%%%%%%%%%%% Slide %%%%%%%%%%%%%%%%%%%%%%%%%%%%%%%%%%%%%%%%%%%%%%%%%%%
%% \heading{Vertex/edge labelling example}
%% 
%% \vspace*{0.2in}
%% 
%% \twofig{\sfile{figures}{edge-labelling13.ps}}{\file{figures}{huffman-clowes.ps}}
%% 
%%%%%%%%%%%% Slide %%%%%%%%%%%%%%%%%%%%%%%%%%%%%%%%%%%%%%%%%%%%%%%%%%%
\heading{Ví dụ về ghi nhãn đỉnh/cạnh}

\vspace*{0.2in}

\twofig{\sfile{figures}{edge-labelling14.ps}}{\file{figures}{huffman-clowes.ps}}

CSP: biến = cạnh, ràng buộc = cấu hình nút có thể

%%%%%%%%%%%% Slide %%%%%%%%%%%%%%%%%%%%%%%%%%%%%%%%%%%%%%%%%%%%%%%%%%%
\heading{Nhận dạng đối tượng}

Ý tưởng đơn giản:\al
-- trích xuất hình dạng 3-D từ hình ảnh\al
-- đối chiếu với ``thư viện hình dạng''

Sự cố:\al
-- trích xuất các bề mặt cong từ hình ảnh\al
-- thể hiện hình dạng của đối tượng được trích xuất\al
-- biểu diễn hình dạng và tính biến đổi của các lớp đối tượng thư viện\al
-- phân đoạn, tắc nghẽn không đúng\al
-- độ sáng, bóng tối, dấu hiệu, tiếng ồn, độ phức tạp không xác định, v.v.

Phương pháp tiếp cận:\al
-- lập chỉ mục vào thư viện bằng cách đo các thuộc tính bất biến của đối tượng\al
-- căn chỉnh tính năng hình ảnh với tính năng đối tượng thư viện được chiếu \al
-- khớp hình ảnh với nhiều chế độ xem được lưu trữ (\emph{khía cạnh}) của đối tượng thư viện\al
-- phương pháp học máy dựa trên số liệu thống kê hình ảnh


%%%%%%%%%%%% Slide %%%%%%%%%%%%%%%%%%%%%%%%%%%%%%%%%%%%%%%%%%%%%%%%%%%
\heading{Nhận dạng chữ số viết tay}

\noindent\framebox[\textwidth]{%
\begin{minipage}{\maxfigwidth%
\noindent{\fboxsep=0,4pt
\framebox{\epsfflex{0.095}{figures/easy21c.eps}}\hfill
\framebox{\epsfflex{0.095}{figures/easy23c.eps}}\hfill
\framebox{\epsfflex{0.095}{figures/easy16c.eps}}\hfill
\framebox{\epsfflex{0.095}{figures/easy12c.eps}}\hfill
\framebox{\epsfflex{0.095}{figures/easy20c.eps}}\hfill
\framebox{\epsfflex{0.095}{figures/easy35c.eps}}\hfill
\framebox{\epsfflex{0.095}{figures/easy13c.eps}}\hfill
\framebox{\epsfflex{0.095}{figures/easy29c.eps}}\hfill
\framebox{\epsfflex{0.095}{figures/easy17c.eps}}\hfill
\framebox{\epsfflex{0.095}{figures/easy19c.eps}\\[6pt]
\framebox{\epsfflex{0.095}{figures/hard53c.eps}}\hfill
\framebox{\epsfflex{0.095}{figures/hard04c.eps}}\hfill
\framebox{\epsfflex{0.095}{figures/hard20c.eps}}\hfill
\framebox{\epsfflex{0.095}{figures/hard13c.eps}}\hfill
\framebox{\epsfflex{0.095}{figures/hard01c.eps}}\hfill
\framebox{\epsfflex{0.095}{figures/hard10c.eps}}\hfill
\framebox{\epsfflex{0.095}{figures/hard34c.eps}}\hfill
\framebox{\epsfflex{0.095}{figures/hard52c.eps}}\hfill
\framebox{\epsfflex{0.095}{figures/hard05c.eps}}\hfill
\framebox{\epsfflex{0.095}{figures/hard14c.eps}}}
\end{minipage}}

3-hàng xóm gần nhất = lỗi 2,4\% \\
400--300--10 đơn vị MLP = lỗi 1,6\%\\
LeNet: 768--192--30--10 đơn vị MLP = lỗi 0,9\%

%%%%%%%%%%%% Slide %%%%%%%%%%%%%%%%%%%%%%%%%%%%%%%%%%%%%%%%%%%%%%%%%%%
\heading{Khớp hình dạng-ngữ cảnh}

Ý tưởng cơ bản: chuyển đổi \emph{hình dạng} (khái niệm quan hệ) thành \\
một tập hợp cố định các thuộc tính \emph{} sử dụng bối cảnh không gian \defn{}\\
của mỗi tập hợp điểm cố định trên bề mặt của hình.

\threefig{figures/exampleA3.eps}{figures/exampleA4.eps}{figures/bin_grid.eps}


%%%%%%%%%%%% Slide %%%%%%%%%%%%%%%%%%%%%%%%%%%%%%%%%%%%%%%%%%%%%%%%%%%
\heading{Tiếp theo là khớp hình dạng-ngữ cảnh.}

Mỗi điểm được mô tả bằng biểu đồ ngữ cảnh cục bộ của nó\\
(số điểm rơi vào mỗi thùng lưới log-cực)

\threefig{figures/exampleA6.eps}{figures/exampleA8.eps}{figures/exampleA7.eps}

%%%%%%%%%%%% Slide %%%%%%%%%%%%%%%%%%%%%%%%%%%%%%%%%%%%%%%%%%%%%%%%%%%
\heading{Tiếp theo là khớp hình dạng-ngữ cảnh.}

Xác định tổng khoảng cách giữa các hình bằng tổng khoảng cách cho
điểm tương ứng dưới sự phù hợp tốt nhất

\epsfxsize=0,3\textwidth
\fig{figures/exampleA5.eps}

Việc học lân cận gần nhất đơn giản mang lại tỷ lệ lỗi 0,63\% trên dữ liệu chữ số NIST


%%%%%%%%%%%% Slide %%%%%%%%%%%%%%%%%%%%%%%%%%%%%%%%%%%%%%%%%%%%%%%%%%%
\heading{Tóm tắt}

Tầm nhìn khó khăn --- nhiễu, mơ hồ, phức tạp

Kiến thức trước là cần thiết để hạn chế vấn đề

Cần kết hợp nhiều tín hiệu: chuyển động, đường viền, tạo bóng, kết cấu, âm thanh nổi

Biểu diễn đối tượng ``Thư viện'': hình dạng và các khía cạnh

So khớp hình ảnh/đối tượng: tính năng, đường nét, vùng, v.v.




\end{huge}
\end{document}


